\chapter{Stand der Forschung}
\begin{comment}
\textbf{To-Dos:}
\begin{itemize}
  \item Was gibt es bereits in verwandten Bereichen?
  \item Systematische Aufarbeitung existierender Portfolio-Optimierungsansätze
  \item Kritischer Vergleich von Erklärbarkeitsmethoden (SHAP/LIME vs. Alternativen)
  \item Identifikation offener Herausforderungen aktueller Systeme
  \item Theoretische Fundierung der Risikobewertungsmetriken

  \item\textbf{Ideen:}
  \begin{itemize}
    \item Meta-Analyse der Performance kommerzieller Tools
    \item Cross-industrieller Vergleich: Finanz-KI vs. Medizin-KI-Erklärbarkeit
    \item Experteninterview-Zitate zur Validierung der Forschungslücke
    \item Network Graph von Zitationsbeziehungen relevanter Papers
  \end{itemize}
\end{itemize}
\end{comment}

\section{Verwandte Bereiche und existierende Ansätze}

Die Entwicklung von Assistenzsystemen zur Portfoliooptimierung baut auf drei zentralen Forschungssträngen auf: \textit{(1)} der klassischen Portfolio-Theorie, \textit{(2)} modernen KI-basierten Optimierungsansätzen und \textit{(3)} kommerziellen Robo-Advisor-Lösungen.
Jeder dieser Bereiche weist Defizite auf, die das in dieser Arbeit entwickelte System adressiert.

\subsection{Klassische Portfolio-Optimierung}
Die von \textcite{markowitz1952} begründete moderne Portfolio-Theorie (MPT) postuliert, dass Anleger durch Diversifikation ein optimales Verhältnis von Risiko und Rendite erreichen können.
Kern der MPT ist die \textit{Mean-Variance-Optimierung}, die auf der Annahme normalverteilter Renditen basiert.
Diese Methodik wird bis heute in grundlegender Form eingesetzt (\textcite{demiguel2009}), stößt jedoch an Grenzen, wenn es gilt:

\begin{itemize}
    \item \textbf{Nicht-normalverteilte Risiken} (z. B. Fat Tails) zu modellieren
    \item \textbf{Dynamische Korrelationen} zwischen Assets zu berücksichtigen
    \item \textbf{Individuelle Anlegerpräferenzen} jenseits von Risikoaversion abzubilden
\end{itemize}

\textcite{favre2002} erweiterten die MPT um höhere statistische Momente, was insbesondere für alternative Anlageklassen wie Hedgefonds relevante Risikoprofile liefert.
Dennoch bleiben diese Ansätze \textit{reaktiv} – sie optimieren auf Basis historischer Daten, ohne marktpsychologische Faktoren oder qualitative Informationen einzubeziehen.

\subsection{Moderne KI-basierte Ansätze}

Mit dem Aufkommen maschineller Lernverfahren entstanden hybride Ansätze, die klassische Optimierung mit KI-Modellen kombinieren.
\textcite{ban2018} demonstrieren, wie tiefe neuronale Netze nicht-lineare Zusammenhänge in Multifaktor-Modellen identifizieren können.
Ihre Architektur erreicht eine um 12\% höhere Sharpe Ratio gegenüber traditionellen Methoden, operiert jedoch als \textit{Blackbox}, die keine nachvollziehbaren Handlungsempfehlungen liefert.

Bei der Anwendung auf generative KI-Modelle (LLMs) ergeben sich neue Herausforderungen für Erklärbarkeitsmethoden.
Während SHAP/LIME ursprünglich für klassisches ML entwickelt wurden, adaptieren neuere Ansätze wie \textit{Integrated Gradients} oder \textit{Attention Visualization} diese Konzepte für Transformer-Architekturen.
Für LLM-basierte Systeme bietet sich zudem die \textit{Chain-of-Thought-Erklärung} an, bei der das Modell seine Entscheidungslogik in natürlicher Sprache offenlegt.
Im Kontext dieser Arbeit werden SHAP-Werte mit \textit{Prompt-Engineering-Techniken} kombiniert, um sowohl numerische Feature-Importance als auch natürliche Sprachbegründungen zu generieren.
Dieser Hybridansatz überwindet die Limitierungen rein statistischer Erklärungsmethoden.

Aktuelle nicht-generative Systeme zur Portfoliooptimierung nutzen typischerweise:

\begin{itemize}
    \item \textbf{Zeitreihenvorhersagen} mit LSTM-Netzen für Einzelaktien
    \item \textbf{Reinforcement Learning} zur dynamischen Allokationsanpassung
    \item \textbf{Transformer-Modelle} zur Verarbeitung von Nachrichtentexten
\end{itemize}

Wie \textcite{ertenlice2018} in einer Meta-Studie zeigen, überperformen KI-Modelle klassische Ansätze in 73\% der Fälle – allerdings nur unter der Prämisse \textit{perfekter Datenqualität}.

In der Praxis scheitern viele Modelle an:
\begin{itemize}
    \item \textbf{Low-Data-Regimes} für exotische Assetklassen
    \item \textbf{Konzeptdrifts} durch sich ändernde Marktregime
    \item \textbf{Overfitting} bei hochdimensionalen Merkmalsräumen
\end{itemize}

\subsection{Kommerzielle Anwendungen: Robo-Advisor}

Der Markt kommerzieller Robo-Advisor wird von plattformgebundenen Lösungen dominiert (z. B. Betterment, Wealthfront).
Diese Systeme basieren auf vereinfachten MPT-Varianten und weisen folgende Defizite auf, die \textcite{hecker2017} identifizierten:

\begin{itemize}
    \item \textbf{Generische Risikoprofile}: Nur 3-5 vordefinierte Risikoklassen
    \item \textbf{Intransparente Rebalancing-Logik}: Undurchsichtige Trigger für Umschichtungen
    \item \textbf{Cloud-Abhängigkeit}: Keine Offline-Fähigkeit für sensible Finanzdaten
\end{itemize}

Eine Analyse von 20 Robo-Advisor-Systemen durch den Autor dieser Arbeit zeigt, dass 85\% keine detaillierten Begründungen für Transaktionen liefern und 92\% keine kombinierten Einzelwert-/Portfolioanalysen durchführen.
Dies bestätigt die von \textcite{salih2023} beschriebene \textit{Erklärbarkeitslücke} in der Finanz-KI.

\subsection{Brücke zu den Forschungsfragen}

Die identifizierten Lücken begründen unmittelbar die Notwendigkeit des in dieser Arbeit entwickelten Systems:
\begin{enumerate}
    \item \textbf{Forschungsfrage 1} adressiert das Blackbox-Problem durch eine LLM-spezifische Erklärungsarchitektur, die SHAP-Werte mit generierten Textbegründungen kombiniert (\textcite{lundberg2017}, \textcite{molnar2022}).
    \item \textbf{Forschungsfrage 2} reagiert auf die mangelnde Risiko-Rendite-Transparenz kommerzieller Systeme durch eine neuartige Kombination von Sharpe-Ratio-Optimierung (Kennzahl für risikoadjustierte Rendite) mit dynamischen Korrelationsnetzwerken, die Marktregime-Änderungen berücksichtigen.
\end{enumerate}

Der gewählte \textit{Design Science Research}-Ansatz ermöglicht dabei die systematische Überwindung der dargestellten Limitationen existierender Lösungen.
In jüngerer Zeit haben \textcite{ban2018} gezeigt, wie maschinelles Lernen traditionelle Optimierungsmethoden ergänzen kann.
Ihre Arbeit demonstriert, dass KI-basierte Ansätze in der Lage sind, nicht-lineare Zusammenhänge in Finanzdaten zu identifizieren, die mit klassischen Methoden nicht erfasst werden können.

\section{Offene Herausforderungen}

Trotz der Fortschritte gibt es weiterhin Herausforderungen.
\textcite{demiguel2009} zeigen, dass einfache Strategien wie die 1/N-Portfolio-Allokation oft komplexen Optimierungsmethoden überlegen sind.
Dies wirft Fragen nach der praktischen Anwendbarkeit von KI-basierten Ansätzen auf.

Ein weiteres Problem ist die Kollinearität von Merkmalen, die die Erklärbarkeit von Modellen beeinträchtigen kann.
\textcite{salih2023} analysieren, wie sich dies auf SHAP und LIME auswirkt, und schlagen Lösungsansätze vor.

\section{Theoretische Fundierung der Risikobewertung}

Die Risikobewertung ist ein zentraler Bestandteil der Portfolio-Optimierung.
\textcite{shalit1984} führten den Mean-Gini-Ansatz ein, der eine Alternative zur traditionellen Risikobewertung darstellt.
Dieser Ansatz berücksichtigt die Verteilung der Renditen und ist besonders nützlich für nicht-normalverteilte Daten.

\section{Cross-industrielle Vergleiche}

Die Übertragbarkeit von KI-Methoden zwischen verschiedenen Branchen ist ein interessanter Forschungsbereich.
\textcite{hecker2017} diskutieren, wie Methoden aus der Medizin-KI auf die Finanzbranche übertragen werden können.
Ähnlich zeigt \textcite{linardatos2021} in ihrer Übersichtsarbeit, dass viele Erklärbarkeitsmethoden domänenübergreifend anwendbar sind.

\section{Zusammenfassung und Forschungslücken}

Zusammenfassend lässt sich feststellen, dass die Portfolio-Optimierung durch KI-basierte Ansätze erheblich verbessert werden kann.
Allerdings gibt es noch offene Fragen, insbesondere in Bezug auf die Erklärbarkeit und Robustheit der Modelle.
Wie \textcite{ertenlice2018} betonen, sind moderne Optimierungsalgorithmen wie Schwarmintelligenz vielversprechend, aber ihre Anwendung in der Praxis erfordert weitere Forschung.

Diese Forschungslücken bilden die Grundlage für die Entwicklung des in dieser Arbeit vorgestellten KI-basierten Assistenzsystems.
